\section{Modelo, resultados y an\'alisis}

% Aviso a\~nadido el 2026-07-27. Las cifras del protocolo con fuga se retiraron de las
% tablas de esta secci\'on el mismo d\'ia (ver comentarios en cada tabla); queda el aviso
% porque la narrativa de la secci\'on sigue siendo la de la defensa y conviene decirlo.
% El t\'itulo va en blanco, no en msorange: naranja sobre el azul de la barra da 2.58:1,
% por debajo del 4.5:1 exigible para texto. El \'enfasis lo carga el contenido.
\begin{frame}{Aviso antes de esta secci\'on}
  \vspace{0.6em}
  \begin{block}{Las tablas de esta secci\'on ya son las honestas}
    \small
    Esta secci\'on presentaba originalmente los AUROC del \textbf{protocolo con fuga} como
    logros del modelo. \textbf{Esas cifras se retiraron:} las tablas que siguen contienen
    las mediciones bajo evaluaci\'on correcta, sobre el grafo corregido.

    \vspace{0.4em}
    Lo que se conserva de la versi\'on de la defensa es la \textbf{narrativa} --- el orden
    en que se plante\'o el problema y se compararon arquitecturas --- porque el art\'iculo
    que salga de aqu\'i trata precisamente de c\'omo se produce y se cree un n\'umero
    inflado.
  \end{block}

  \vspace{0.3em}
  \begin{block}{D\'onde est\'a cada cosa}
    \scriptsize
    \begin{itemize}
      \item \textbf{Secci\'on 4} --- los cuatro experimentos que desarmaron el resultado
            original, y la atribuci\'on de la ca\'ida a sus dos causas.
      \item \textbf{Secci\'on 5} --- el estado actual: grafo corregido, 4 semillas, y la
            evidencia de que la inflaci\'on es \textbf{del protocolo} y no de este modelo.
            \emph{Esa secci\'on tiene la \'ultima palabra sobre cualquier cifra de aqu\'i.}
    \end{itemize}
  \end{block}
\end{frame}

\begin{frame}{Arquitectura del modelo}
  \begin{columns}[T]
    \begin{column}{0.56\textwidth}
      \begin{block}{Modelo principal}
        `miRNAGraphTransformer` con capas HGT sobre un grafo heterog\'eneo de 3 tipos de nodos y 5 tipos de aristas.
      \end{block}

      \small
      \begin{itemize}
        \item 4 capas HGT en la mejor variante V2.
        \item 8 heads de atenci\'on.
        \item 512 canales ocultos.
        \item Dos cabezas de salida:
        \begin{itemize}
          \item predicci\'on de enlaces miRNA-gen,
          \item clasificaci\'on de tipo celular.
        \end{itemize}
        \item Funci\'on de p\'erdida combinada: $L = L_{link} + 0.5L_{clf} + 0.01L_{sparsity}$.
      \end{itemize}
    \end{column}
    \begin{column}{0.40\textwidth}
      \begin{tikzpicture}[every node/.style={font=\small}]
        \node[draw=msblue, rounded corners, fill=mslight, minimum width=3.9cm, minimum height=0.9cm, align=center] (i) at (0, 3.4) {HeteroData};
        \node[draw=msblue, rounded corners, fill=mslight, minimum width=3.9cm, minimum height=0.9cm, align=center] (h1) at (0, 2.1) {HGT Layer 1};
        \node[draw=msblue, rounded corners, fill=mslight, minimum width=3.9cm, minimum height=0.9cm, align=center] (h2) at (0, 0.9) {HGT Layer 2};
        \node[draw=msblue, rounded corners, fill=mslight, minimum width=3.9cm, minimum height=0.9cm, align=center] (h3) at (0, -0.3) {HGT Layer 3};
        \node[draw=msblue, rounded corners, fill=mslight, minimum width=3.9cm, minimum height=0.9cm, align=center] (h4) at (0, -1.5) {HGT Layer 4};
        \node[draw=msgreen, rounded corners, fill=green!10, minimum width=1.8cm, minimum height=0.85cm, align=center] (o1) at (-1.2, -3.1) {Link\\head};
        \node[draw=msorange, rounded corners, fill=orange!10, minimum width=1.8cm, minimum height=0.85cm, align=center] (o2) at (1.2, -3.1) {Cell type\\head};
        \draw[->, thick] (i) -- (h1);
        \draw[->, thick] (h1) -- (h2);
        \draw[->, thick] (h2) -- (h3);
        \draw[->, thick] (h3) -- (h4);
        \draw[->, thick] (h4) -- (o1);
        \draw[->, thick] (h4) -- (o2);
      \end{tikzpicture}
    \end{column}
  \end{columns}
\end{frame}

\begin{frame}{Detalle t\'ecnico: mec\'anica interna de las capas HGT}
  \small
  \begin{block}{Flujo de c\'omputo}
    \scriptsize (`models/layers.py`, `models/hetero\_gnn.py`) \small
    \begin{enumerate}
      \item \textbf{NodeProjector}: `Linear(-1, hidden)` + ReLU por tipo de nodo (proyecci\'on \emph{lazy} de PyG: se adapta a la dimensi\'on real de entrada en el primer forward, sin fijarla a mano).
      \item \textbf{HGTLayer $\times$ L}: envuelve `torch\_geometric.nn.HGTConv` (Hu et al., WWW 2020) --- atenci\'on espec\'ifica por tipo de relaci\'on, seguida de \textbf{conexi\'on residual} + \textbf{LayerNorm} por tipo de nodo, y dropout tras la convoluci\'on.
      \item \textbf{TargetPredictor}: concatena `[emb\_miRNA | emb\_gene]` ($2\times$hidden) $\to$ MLP (Linear--ReLU--Dropout 0.1--Linear) $\to$ logit escalar.
      \item \textbf{CellTypeClassifier}: MLP (hidden $\to$ hidden/2 $\to$ 11 clases), dropout 0.1.
    \end{enumerate}
  \end{block}

  \begin{itemize}
    \item Las dos cabezas comparten el mismo backbone de embeddings --- no hay ramas separadas antes de la \'ultima capa HGT.
    \item El t\'ermino de sparsity act\'ua sobre el embedding de miRNA que sale de este backbone.
  \end{itemize}
\end{frame}

\begin{frame}{Detalle t\'ecnico: protocolo de entrenamiento y muestreo}
  \small
  \begin{itemize}
    \item \textbf{Mini-batching}: el grafo completo vive en CPU; `NeighborLoader` muestrea vecindarios por batch (`num\_neighbors=[10, 5]`, uno por capa HGT; `batch\_size=512`).
    \item \textbf{Negative sampling}: \emph{(versi\'on original)} un negativo \emph{uniforme} al azar por cada arista positiva. \textbf{Corregido:} ahora negativos \textbf{pareados por grado} (`splits.LinkSampler`) --- ver Experimento 2.
    \item \textbf{Partici\'on de datos}: split 80/10/10 \textbf{sobre nodos `cell`}. \emph{(versi\'on original:} no hab\'ia partici\'on a nivel de arista miRNA-gen\emph{)}. \textbf{Corregido:} holdout expl\'icito de aristas (`splits.build\_edge\_split`, `RandomLinkSplit` con `rev\_edge\_types`).
    \item \textbf{Optimizaci\'on}: AdamW + `CosineAnnealingLR` (`T\_max = num\_epochs`, 200 en V2); gradient clipping (`max\_norm=1.0`); early stopping por `val\_loss` (patience configurable).
    \item \textbf{Multi-GPU}: DDP con backend NCCL, `torchrun --nproc\_per\_node=4` sobre 4$\times$ A100; desbalance de clases en link prediction compensable v\'ia `pos\_weight` en la BCE (`CombinedLoss`, actualmente no fijado en los configs revisados).
  \end{itemize}
\end{frame}

\begin{frame}{Configuraci\'on V2: variante ganadora}
  \small
  \begin{tabularx}{\textwidth}{>{\bfseries}p{3.2cm}X}
    \toprule
    Par\'ametro & Valor \\
    \midrule
    hidden\_channels & 512 \\
    num\_layers & 4 \\
    dropout & 0.3 \\
    learning rate & $5 \times 10^{-4}$ \\
    batch\_size & 512 \\
    num\_epochs & 200 \\
    patience & 25 \\
    vecindad de muestreo & [10, 5] \\
    hardware & 4x NVIDIA A100-SXM4-80GB \\
    optimizador & Adam, weight decay $10^{-5}$ \\
    \bottomrule
  \end{tabularx}

  \vspace{0.5em}
  \begin{itemize}
    \item Esta configuraci\'on profundiza y ensancha el baseline, a la vez que baja el learning rate y aumenta la regularizaci\'on.
    \item La evidencia del repositorio y del job de comparaci\'on 5042 apunta a que V2 es la mejor combinaci\'on estable encontrada hasta ahora para el objetivo completo.
  \end{itemize}
\end{frame}

% Tabla reemplazada el 2026-07-27. Antes conten\'ia los AUROC del protocolo con fuga
% (HGT 0.9836, ablaci\'on 0.9374, GCN 0.9170...), presentados como logros. Ahora contiene
% los del protocolo honesto medidos el mismo d\'ia sobre el grafo corregido (job 5849,
% results/comparison/comparison_table_checkpoints_v3fixed_baselines_edgesplit.tsv).
% Cambiar la tabla invierte su conclusi\'on, y eso es correcto: el GCN simple le gana.
\begin{frame}{Comparativa entre modelos --- protocolo honesto}
  \small
  \begin{tabular}{lcccc}
    \toprule
    \textbf{Modelo} & \textbf{AUROC} & \textbf{AUPRC} & \textbf{acc} & \textbf{F1} \\
    \midrule
    \textbf{GCN homog\'eneo} & \textbf{0.6236} & \textbf{0.6466} & 0.8498 & 0.8114 \\
    HGT sin co-expresi\'on & 0.6185 & 0.6452 & 0.9675 & 0.9468 \\
    HGT V2 \emph{(el modelo del proyecto)} & 0.6080 & 0.6371 & 0.9565 & 0.9301 \\
    MLP \emph{(sin grafo)} & 0.5202 & 0.5202 & 0.8948 & 0.8559 \\
    Random \emph{(control)} & 0.4938 & 0.4932 & 0.0891 & 0.0784 \\
    Ablation no miRNA & --- & --- & \textbf{0.9923} & \textbf{0.9868} \\
    \bottomrule
  \end{tabular}

  \vspace{0.4em}
  \begin{block}{Lo que dice esta tabla}
    \scriptsize
    Aristas retenidas y controles pareados por popularidad, todos los modelos en
    condiciones id\'enticas. \textbf{El GCN homog\'eneo (0.6236) le gana al HGT (0.6080)}:
    la arquitectura heterog\'enea no est\'a pagando su costo. El MLP sin grafo (0.5202) no
    llega ni a la heur\'istica Adamic--Adar (0.5912), y `random' cae en 0.4938 --- justo
    donde debe caer el azar, que es la prueba de que la m\'etrica se calcula bien.

    \vspace{0.2em}
    \textbf{`no\_miRNA' no tiene AUROC porque se le quitaron esas aristas:} la tarea no
    existe para \'el. Y aun as\'i tiene la mejor exactitud celular (0.9923) --- por eso
    comparar p\'erdidas totales entre modelos con objetivos distintos no significa nada.
  \end{block}
\end{frame}

\begin{frame}{C\'omo defender la comparativa (versi\'on corregida)}
  \begin{columns}[T]
    \begin{column}{0.5\textwidth}
      % Este bloque afirmaba que la ventaja del HGT en regulaci\'on era ``real pero
      % moderada''. Era falso por dos v\'ias: el n\'umero que la sosten\'ia era inflaci\'on
      % del protocolo, y bajo evaluaci\'on honesta el GCN simple le gana. Reescrito.
      \begin{block}{Lo que la tabla \emph{s\'i} sostiene}
        \small
        Que la \textbf{clasificaci\'on de tipo celular es s\'olida} en todos los modelos que
        la hacen (0.85--0.99), y que el control aleatorio cae en el azar --- se\~nal de que
        la m\'etrica est\'a bien calculada.
      \end{block}
      \begin{block}{Lo que \textbf{no} sostiene}
        \small
        \textbf{No sostiene que el HGT sea el mejor en regulaci\'on: no lo es.} Bajo
        evaluaci\'on honesta el GCN homog\'eneo (0.6236) le gana al HGT (0.6080), y ninguno
        de los dos se despega de una heur\'istica de 1999 (0.5912). \textbf{Reclamar
        superioridad arquitect\'onica aqu\'i es invitar a que un revisor lo desmonte.}
      \end{block}
    \end{column}
    \begin{column}{0.47\textwidth}
      \begin{block}{Las dos tareas se comportan distinto}
        \small
        \begin{itemize}
          \item \textbf{Clasificaci\'on celular}: casi todos los modelos superan 0.89. La identidad celular est\'a \emph{muy} marcada en los datos; distingue poco entre arquitecturas.
          \item \textbf{Link prediction}: es la tarea cercana a mecanismo, y la que \emph{deber\'ia} discriminar entre modelos. Bajo evaluaci\'on honesta apenas lo hace: todos caen entre 0.49 y 0.62.
        \end{itemize}
      \end{block}
      \vspace{0.2em}
      \begin{block}{Mensaje para defensa}
        \scriptsize
        ``La ventaja de V2 est\'a en la tarea regulatoria, no en la clasificatoria --- y la reportamos con su magnitud real, no inflada.''
      \end{block}
    \end{column}
  \end{columns}
\end{frame}

\begin{frame}{Separar dos tareas: clasificar c\'elulas vs inferir regulaci\'on}
  \begin{columns}[T]
    \begin{column}{0.48\textwidth}
      \begin{block}{Tarea 1: clasificaci\'on celular}
        \small
        Pregunta: \textquestiondown{}a qu\'e tipo celular pertenece esta observaci\'on?
      \end{block}
      \begin{itemize}
        \item M\'etricas: `cell_acc`, `cell_f1`
        \item Problema m\'as f\'acil
        \item Puede seguir alto en ablaciones
        \item Valida la representaci\'on celular del modelo
      \end{itemize}
    \end{column}
    \begin{column}{0.48\textwidth}
      \begin{block}{Tarea 2: predicci\'on regulatoria}
        \small
        Pregunta: \textquestiondown{}qu\'e par miRNA-gen es plausible en este contexto biol\'ogico?
      \end{block}
      \begin{itemize}
        \item M\'etricas: AUROC, AUPRC
        \item Problema m\'as dif\'icil y m\'as cercano a mecanismo
        \item Es la tarea que \emph{deber\'ia} discriminar entre arquitecturas --- y honestamente casi no lo hace (0.49 a 0.62)
        \item \textbf{No} es el valor diferencial del HGT: honestamente el GCN simple le gana (0.6236 vs 0.6080)
      \end{itemize}
    \end{column}
  \end{columns}

  \vspace{0.5em}
  \begin{block}{Mensaje que conviene decir de forma expl\'icita}
    \small
    Una ablaci\'on puede retener excelente accuracy celular y aun as\'i perder valor biol\'ogico si ya no modela bien la regulaci\'on miRNA$\rightarrow$gen.
  \end{block}
\end{frame}

\begin{frame}{M\'etricas en test --- y por qu\'e hay que leerlas con cuidado}
  \begin{columns}[T]
    \begin{column}{0.52\textwidth}
      \small
      \begin{tabular}{lcc}
        \toprule
        \textbf{M\'etrica} & \textbf{Protocolo} & \textbf{Protocolo} \\
                           & \textbf{original} & \textbf{honesto} \\
        \midrule
        % El protocolo original se cita ahora con su medici\'on multi-semilla sobre el
        % grafo corregido (0.9867 +/- 0.0011, n=4) en lugar del 0.9836 de una sola semilla
        % sobre el grafo con errores: mismo protocolo, medido bien y reproducible.
        AUROC   & 0.9867 & \textbf{0.6276} \\
        AUPRC   & 0.9850 & \textbf{0.6598} \\
        % La clasificaci\'on celular no se ve afectada por la fuga de aristas, as\'i que no
        % tiene dos versiones: se reporta una sola vez, la de prueba.
        Cell acc & \multicolumn{2}{c}{\textbf{0.9916} \emph{(no afectada)}} \\
        \bottomrule
      \end{tabular}

      \vspace{0.6em}
      \begin{itemize}
        \scriptsize
        \item La clasificaci\'on celular \textbf{no cambia}: es real bajo ambos protocolos.
        \item El AUROC de regulaci\'on \textbf{s\'i} cambia, y mucho. Ver Experimentos 3 y 4.
      \end{itemize}
    \end{column}
    \begin{column}{0.45\textwidth}
      \begin{block}{Advertencia}
        \scriptsize
        La columna ``original'' mide qu\'e tan bien el modelo \textbf{reconstruye} interacciones que ya vio, contra controles no pareados. \textbf{No mide predicci\'on.}

        \vspace{0.2em}
        Las diapositivas siguientes (saliencia, circuitos, enriquecimiento) derivan \emph{todas} del modelo evaluado con el protocolo original. Se presentan como \textbf{hip\'otesis preliminares}, no como hallazgos.
      \end{block}
    \end{column}
  \end{columns}
\end{frame}

% ---------------------------------------------------------------------------
% Bloque retirado el 2026-07-27. Aqui habia 7 marcos de interpretabilidad:
%   - 3 figuras de la PRIMERA version (mayo 2026): 01_umap_embeddings,
%     02_mirna_heatmap y 05_auroc_auprc. Las dos ultimas son insostenibles:
%     05_auroc_auprc son las curvas ROC del modelo CON FUGA, con pie que hablaba
%     de "la robustez del modelo"; 02_mirna_heatmap es saliencia de features que
%     la auditoria probo reproducibles con torch.randn (max|diff| = 0.0), es decir
%     sin contenido -- una imagen de ruido con pie biologico.
%   - "Hallazgos interpretables" y "Circuitos reguladores de alta confianza"
%     (scores 0.9699, 0.9558...): son los no-hallazgos suspendidos, ordenados por
%     un cabezal que supera a una heuristica sin aprendizaje por 3.6 puntos, y que
%     ademas no puede ser especifico de tipo celular (ver marco siguiente).
%   - Un marco hablaba de "reguladores transversales del fenotipo EM", claim que
%     la decision D1 retira.
% Si algun dia se quiere una figura del espacio celular (la clasificacion celular
% SI es real), se regenera desde un checkpoint v3fixed; no se reutiliza la de mayo.
% ---------------------------------------------------------------------------

\begin{frame}{Por qu\'e no hay diapositivas de circuitos}
  \begin{block}{No es prudencia: es que la arquitectura no pod\'ia producirlos}
    \small
    `TargetPredictor' recibe \textbf{solo} el miRNA y el gen. \textbf{No recibe la
    c\'elula.} Y el an\'alisis puntuaba cada par miRNA$\to$gen \textbf{una sola vez, de
    forma global}, para luego filtrar esa lista \'unica por los miRNAs m\'as salientes de
    cada tipo celular.
  \end{block}

  \vspace{0.3em}
  \onslide<2->{%
  \begin{block}{La consecuencia}
    \small
    \textbf{`miR-23a-3p $\to$ CCL7' ten\'ia exactamente el mismo score en todos los tipos
    celulares.} La ``especificidad de tipo celular'' viv\'ia \emph{solo} en el filtro de
    saliencia, nunca en el score de la interacci\'on. La premisa central ---
    regulaci\'on espec\'ifica de tipo celular --- \textbf{no lleg\'o a implementarse}, y con
    esta arquitectura no pod\'ia.
  \end{block}}

  \vspace{0.3em}
  \onslide<3->{%
  \begin{block}{Y esto s\'i es un aporte, no solo una limitaci\'on}
    \scriptsize
    Es \textbf{arquitect\'onico, no estad\'istico}: no lo arregla reentrenar. Y generaliza
    m\'as all\'a de este proyecto --- \emph{toda afirmaci\'on de predicci\'on
    espec\'ifica de contexto deber\'ia comprobarse contra tres preguntas: \textbf{(1)}
    \textquestiondown{}la variable de contexto est\'a en el artefacto? \textbf{(2)}
    \textquestiondown{}la recibe el cabezal que puntua? \textbf{(3)}
    \textquestiondown{}la evaluaci\'on var\'ia con ella?} Nosotros fallamos las tres.
  \end{block}}
\end{frame}